\documentclass{article}
\usepackage[UTF8]{ctex}
\usepackage{amsmath,amssymb,amsfonts}
\usepackage{enumitem}
\usepackage{geometry}
\geometry{a4paper, margin=1in}
% === 错题本排版优化（不参与题目标识） ===
\setlength{\parindent}{0pt}
\setlength{\parskip}{0.35em plus 0.10em minus 0.08em}
\setlength{\emergencystretch}{3em}
\setlist{itemsep=0.25em, topsep=0.40em, parsep=0pt}
\allowdisplaybreaks[2]
% === 错题本排版优化结束 ===
\begin{document}
\section*{南京理工大学2024年高等代数真题}
\section*{89．南京理工大学2024年高等代数试卷}
一．填空题．每题5分，共40分．\\

\subsection{题目1}
1．设 \(f(x)=x^{4}+a, g(x)=x^{6}+b\) ，若 \((f(x), g(x))=x^{2}+1\) ，求 \(a b\) 的值 \(\_\_\_\_\)．

\subsection{题目2}
2．设 \(n\) 阶行列式 \(D_{n}=\left|a_{i j}\right|=1\) ，令 \(b_{i j}=3^{2 i-j} a_{i j}\) ，求 \(\left|b_{i j}\right|\) 的值 \(\_\_\_\_\)．

\subsection{题目3}
3．设 \(A\) 为一个3阶方阵，将 \(A\) 的第1行加到第3行，再将第1行与第2行互换之后为 \(2 E\) ，求矩阵

\subsection{题目4}
4．设 \(A^{*}=\left(\begin{array}{ccc}1 & -1 & 1 \\ a & 2 & -b \\ 2 a & -3 a-2 & 3 a+b\end{array}\right)\) 为矩阵 \(A\) 的伴随矩阵，则 \(a, b\) 依次为 \(\_\_\_\_\)

\subsection{题目5}
5．设 \(A\) 为一个 3 阶方阵，且 \(|A+E|=|A+2 E|=|A+3 E|=0\) ，求 \(\left|A-A^{-1}\right|=\) \(\_\_\_\_\)．\\
\(\_\_\_\_\)

\subsection{题目6}
6．设二次型 \(f\left(x_{1}, x_{2}, x_{3}\right)=X^{\mathrm{T}} A X\) ，其中 \(A^{\mathrm{T}}=A,|A|=a, r(A+b E)=1\) ，若 \(f\) 正定，求 \(a, b\) 满足的条件为 \(\_\_\_\_\)．

\subsection{题目7}
7．在几何空间中，设 \(O-x y z\) 为一直角坐标系， \(\mathcal{A}\) 表示将空间绕 \(O x\) 轴，由 \(O y\) 轴向 \(O z\) 轴旋转 \(90^{\circ}\) 的线性变换，则 \(\mathcal{A}\) 在基 \(\varepsilon_{1}=(1,0,0), \varepsilon_{2}=(0,1,0), \varepsilon_{3}=(0,0,1)\) 下的矩阵为 \(\_\_\_\_\)．\\
\(\_\_\_\_\)

\subsection{题目8}
8．设一数域 \(P\) ，设线性空间 \(P^{2}\) 中的基 \(\alpha_{1}, \alpha_{2}\) 的对偶基为 \(f_{1}, f_{2}\) ，则 \(P^{2}\) 中的基 \(\alpha_{1}+\alpha_{2}, \alpha_{1}-\alpha_{2}\) 的对偶基为 \(\_\_\_\_\)（用 \(f_{1}, f_{2}\) 表示）．

二．（10 分）设实数域上的多项式为 \(f(x)=x^{3}+6 x^{2}+3 p x+8\) ，求当 \(p\) 为何值时，\(f(x)\) 有重因式？

三．（10 分）设 \(A=\left(a_{i j}\right)\) 为一个 \(n\) 阶矩阵，其中 \(a_{i j}=\max \{i, j\}\) ，求 \(\left|A^{*}\right|\) 。

四．（15分）设 \(A=\left(\alpha_{1}, \alpha_{2}, \alpha_{3}, \alpha_{4}\right)\) 为一个 4 阶方阵，且 \(\alpha_{1}, \alpha_{2}\) 线性无关，\(\alpha_{3}=\alpha_{1}+\alpha_{2}, \alpha_{4}=\alpha_{1}-\alpha_{2}\)，求方程组 \(A X=\alpha_{3}+\alpha_{4}\) 的通解。

五。（15 分）设 \(V_{1}\) 为 \(\alpha_{1}=(1,2,3), \alpha_{2}=(1,0,1)\) 生成的子空间，\(V_{2}\) 为 \(\beta_{1}=(-1,2, t), \beta_{2}=(4,1,5)\)生成的子空间．若 \(V_{1}=V_{2}\) ，求 \(t\) 的值，并将 \(\beta_{1}, \beta_{2}\) 写成 \(\alpha_{1}, \alpha_{2}\) 的线性组合．

六．（15 分）设矩阵 \(A=\left(\begin{array}{lll}1 & a & 1 \\ a & 2 a & 1 \\ 1 & 1 & 1\end{array}\right), B=\left(\begin{array}{lll}1 & 0 & 0 \\ 0 & 1 & 0 \\ 0 & 0 & 0\end{array}\right)\) ，且 \(A\) 与 \(B\) 合同．\\
（1）（5 分）求 \(a\) 的值．\\
（2）（10 分）求一可逆矩阵 \(C\) ，使得 \(C^{\mathrm{T}} A C=B\) 。

七．（15 分）设 \(V=P[x]_{n}\) 为次数小于 \(n\) 多项式与零多项式生成的线性空间，在 \(V\) 上定义一个二元函数

\[
\varphi(f(x), g(x))=\int_{-1}^{1} f(x) g(x) \mathrm{d} x, f(x), g(x) \in V
\]

（1）（6 分）若 \(n=4\) ，求 \(\varphi\) 在基 \(1, x, x^{2}, x^{3}\) 下的度量矩阵。\\
（2）（9 分）证明：\(\varphi\) 为非退化的．

八．（15 分）设矩阵 \(A=\left(\begin{array}{ccc}0 & 2 & -2 \\ -1 & 3 & -3 \\ 1 & -2 & a\end{array}\right), B=\left(\begin{array}{ccc}1 & -2 & 0 \\ 0 & b & 0 \\ 0 & 3 & 1\end{array}\right)\) ，且 \(A\) 与 \(B\) 相似．

\begin{itemize}
  \item[（1）] （ 5 分）求 \(a, b\) 的值．
  \item[（2）] （10 分）求可逆矩阵 \(P\) ，使得 \(P^{-1} A P=B\) ．
\end{itemize}

九．（15 分）设 \(A\) 为一个秩为 \(r\) 的 \(n\) 阶方阵，且 \(A^{2}=A\) 。证明：存在一个秩为 \(r\) 的 \(r \times n\) 矩阵 \(B\) 与一个秩为 \(r\) 的 \(n \times r\) 矩阵 \(C\) ，满足 \(A=C B\) ，且 \(B C=E_{r}\)（ \(E_{r}\) 为 \(r\) 阶单位矩阵）。
\end{document}
